\section{Circuit Complexity --- Circuit Eval Correct}
%%% generated by the blueprint pipeline from TCSlib.Complexity.CircuitComplexity.CircuitEvalCorrect
\subsection*{Circuit Complexity --- Circuit Eval Correct: Introduction}

Declarations of the Lean module \texttt{TCSlib.Complexity.CircuitComplexity.CircuitEvalCorrect}.

\subsection*{Circuit Complexity --- Circuit Eval Correct: Results}

\begin{definition}[Valid gate lists]
\lean{BoolCircuit.CircuitEval.GatesOK}
\label{BoolCircuit.CircuitEval.GatesOK}
\leanok
\uses{BoolCircuit.DAGGate}
A list of gates is valid relative to a current length $L$ if, for each gate in order,
its argument list has at most two distinct indices, a NOT gate has exactly one argument,
and every argument index is less than the current length; after each gate, the current
length increases by one. The empty list is valid for every $L$.
\end{definition}

\begin{theorem}[Folded gate evaluation]
\lean{BoolCircuit.CircuitEval.accFin_foldl}
\label{BoolCircuit.CircuitEval.accFin_foldl}
\leanok
\uses{BoolCircuit.CircuitEval.accComb, BoolCircuit.CircuitEval.accFin, BoolCircuit.CircuitEval.accInit, BoolCircuit.DAGGate.eval, BoolCircuit.GateKind}
\difficulty{4}
For every gate kind $k$, argument list $a$, and Boolean value list $v$, folding the
gate’s accumulator update over the arguments, starting from its prescribed initial
accumulator and then applying its final adjustment, gives the value of the gate of kind
$k$ on arguments $a$ with values read from $v$ (using false when an index is out of
range).
\end{theorem}

\begin{theorem}[Acceptance while reading a gate argument]
\lean{BoolCircuit.CircuitEval.absRun_u}
\label{BoolCircuit.CircuitEval.absRun_u}
\leanok
\uses{BoolCircuit.CircuitEval.LA.rej, BoolCircuit.CircuitEval.absOf, BoolCircuit.CircuitEval.absRun, BoolCircuit.CircuitEval.absStep, BoolCircuit.CircuitEval.accComb, BoolCircuit.CircuitEval.lact, BoolCircuit.CircuitEval.markCell, BoolCircuit.GateKind}
\difficulty{5}
Fix a gate kind $k$, counter $c\in\{0,1,2\}$, Boolean accumulator $b$, Boolean value
list $v$, and marked-index list $M$. For every Boolean string $\ell$ and position $j$,
the evaluator started in its argument-reading phase with these data accepts exactly when
there are a natural number $r$ and a Boolean string $t$ such that $\ell$ consists of $r$
true bits followed by a false bit and then $t$, $j+r<|v|$, and $j+r$ is not marked in
$M$. In that case, running from the argument-accumulation phase with counter increased
by one, accumulator combined with the value of $v$ at index $j+r$ (using OR for an OR
gate and AND otherwise), the updated marked list with $j+r$ added at its front, and
position zero, on $t$ also accepts.
\end{theorem}

\begin{theorem}[Acceptance during output-index parsing]
\lean{BoolCircuit.CircuitEval.absRun_o}
\label{BoolCircuit.CircuitEval.absRun_o}
\leanok
\uses{BoolCircuit.CircuitEval.LA.go, BoolCircuit.CircuitEval.LA.rej, BoolCircuit.CircuitEval.absOf, BoolCircuit.CircuitEval.absRun, BoolCircuit.CircuitEval.absStep, BoolCircuit.CircuitEval.lact}
\difficulty{5}
For every Boolean value list $v$ and marked-index list $M$, and for every Boolean string
$\ell$ and natural number $j$, the evaluator started in its output-index-reading phase
at position $j$ accepts on $\ell$ exactly when there is a natural number $r$ such that
$\ell$ consists of $r$ true bits followed by a false bit, $j+r<|v|$, and the value of
$v$ at index $j+r$ is true.
\end{theorem}

\begin{theorem}[Acceptance during gate-argument parsing]
\lean{BoolCircuit.CircuitEval.absRun_a}
\label{BoolCircuit.CircuitEval.absRun_a}
\leanok
\uses{BoolCircuit.CircuitEval.LA.go, BoolCircuit.CircuitEval.LA.rej, BoolCircuit.CircuitEval.absOf, BoolCircuit.CircuitEval.absRun, BoolCircuit.CircuitEval.absRun_u, BoolCircuit.CircuitEval.absStep, BoolCircuit.CircuitEval.accComb, BoolCircuit.CircuitEval.accFin_foldl, BoolCircuit.CircuitEval.accInit, BoolCircuit.CircuitEval.lact, BoolCircuit.DAGGate.eval, BoolCircuit.GateKind, BoolCircuit.encodeList, BoolCircuit.encodeNat}
\difficulty{7}
Fix a gate kind $k$, a list of Boolean values $v$, and natural number $m$. For every
list $p$ of natural numbers, counter $c\in\{0,1,2\}$, marked-index list $M$, and Boolean
string $\ell$, suppose $|p|+m=2$, $c=|p|$, the entries of $p$ are distinct and all lie
below $|v|$, and an index belongs to $M$ exactly when it belongs to $p$. Initialize the
evaluator in its argument-reading phase with counter $c$ and accumulator obtained by
starting with false for an OR gate and true otherwise, then combining the values indexed
by $p$ using OR for an OR gate and AND otherwise. The run accepts exactly when there are
lists $a,r$ such that $\ell$ is the encoding of the list $a$ followed by $r$, the
concatenation $p++a$ has length at most two and has distinct entries all below $|v|$,
and, for a NOT gate, that concatenation has length one; moreover, the evaluator started
in its gate-reading phase with the gate value of kind $k$ on indices $p++a$ appended to
$v$, and run on $r$, accepts.
\end{theorem}

\begin{theorem}[Acceptance during circuit gate parsing]
\lean{BoolCircuit.CircuitEval.absRun_g}
\label{BoolCircuit.CircuitEval.absRun_g}
\leanok
\uses{BoolCircuit.CircuitEval.GatesOK, BoolCircuit.CircuitEval.LA.go, BoolCircuit.CircuitEval.LA.rej, BoolCircuit.CircuitEval.absOf, BoolCircuit.CircuitEval.absRun, BoolCircuit.CircuitEval.absRun_a, BoolCircuit.CircuitEval.absRun_o, BoolCircuit.CircuitEval.absStep, BoolCircuit.CircuitEval.accComb, BoolCircuit.CircuitEval.accInit, BoolCircuit.CircuitEval.kindOf, BoolCircuit.CircuitEval.lact, BoolCircuit.DAGGate.encode, BoolCircuit.DAGGate.eval, BoolCircuit.GateKind, BoolCircuit.GateKind.encode, BoolCircuit.encodeList, BoolCircuit.encodeNat, BoolCircuit.runWith, BoolCircuit.runWith_cons}
\difficulty{7}
For every Boolean string $\ell$ and Boolean value list $v$, the evaluator started in its
gate-parsing phase with values $v$ accepts on $\ell$ exactly when there are a list $G$
of gates and a natural number $o$ such that $\ell$ is the encoding of $G$ followed by
the encoding of $o$, every gate in $G$ has at most two distinct arguments, each NOT gate
has exactly one argument, and every argument of each gate is below the number of values
available when that gate is evaluated. In addition, $o<|v|+|G|$, and the value at index
$o$ after evaluating the gates of $G$ successively from the initial values $v$ is true.
\end{theorem}

\begin{theorem}[Recursive gate conditions and indexed conditions]
\lean{BoolCircuit.CircuitEval.gatesOK_iff}
\label{BoolCircuit.CircuitEval.gatesOK_iff}
\leanok
\uses{BoolCircuit.CircuitEval.GatesOK, BoolCircuit.DAGGate}
\difficulty{4}
For any finite list $G$ of Boolean circuit gates and any natural number $	ext{len}$, the
recursive condition that each gate has at most two distinct arguments, every NOT gate
has exactly one argument, and each argument of the current gate is less than the current
available length (which increases by one at each gate) is equivalent to the conjunction
of these local gate conditions for every gate in $G$ and the indexed condition that the
argument at position $i$ of $G$ is less than $	ext{len}+i$ for every valid index $i$.
\end{theorem}

\begin{theorem}[Leading ones in a padded word]
\lean{BoolCircuit.CircuitEval.leadOnes_replicate}
\label{BoolCircuit.CircuitEval.leadOnes_replicate}
\leanok
\uses{BoolCircuit.CircuitEval.leadOnes}
\difficulty{4}
For every natural number $n$ and every Boolean word $r$, the number of initial
consecutive $1$'s in the word consisting of $n$ copies of $1$, followed by $0$ and then
$r$, is exactly $n$.
\end{theorem}

\begin{theorem}[First zero after the leading ones]
\lean{BoolCircuit.CircuitEval.eq_replicate_leadOnes}
\label{BoolCircuit.CircuitEval.eq_replicate_leadOnes}
\leanok
\uses{BoolCircuit.CircuitEval.leadOnes}
\difficulty{4}
Let $c$ be a finite Boolean word containing at least one $0$, and let $r$ be the number
of its initial consecutive $1$'s. Then $c$ is the concatenation of $r$ copies of $1$,
followed by $0$, followed by the suffix of $c$ after those first $r+1$ entries.
\end{theorem}

\begin{theorem}[Skipping a unary natural-number code]
\lean{BoolCircuit.CircuitEval.absRun_skipN}
\label{BoolCircuit.CircuitEval.absRun_skipN}
\leanok
\uses{BoolCircuit.CircuitEval.LA.go, BoolCircuit.CircuitEval.absOf, BoolCircuit.CircuitEval.absRun, BoolCircuit.CircuitEval.absStep, BoolCircuit.CircuitEval.lact}
\difficulty{4}
For every natural number $n$, Boolean value list $v$, marked-index list $M$, position
$j$, and Boolean string $r$, running the evaluator from its skip-count phase on $n$ true
bits followed by a false bit and then $r$ has the same result as starting in its
gate-parsing phase at position $j$ and running on $r$.
\end{theorem}

\begin{theorem}[Finite-word reconstruction by indexed values]
\lean{BoolCircuit.CircuitEval.ofFn_getD_eq_take}
\label{BoolCircuit.CircuitEval.ofFn_getD_eq_take}
\leanok
\difficulty{3}
For every Boolean word $x$ and natural number $n\le |x|$, the length-$n$ word whose
entry at index $i<n$ is $x_i$ is the prefix of $x$ of length $n$.
\end{theorem}

\begin{theorem}[Verdict for an encoded circuit]
\lean{BoolCircuit.CircuitEval.verdict_pairEncode}
\label{BoolCircuit.CircuitEval.verdict_pairEncode}
\leanok
\uses{BoolCircuit.CircuitEval.absRun_g, BoolCircuit.CircuitEval.absRun_skipN, BoolCircuit.CircuitEval.eq_replicate_leadOnes, BoolCircuit.CircuitEval.gatesOK_iff, BoolCircuit.CircuitEval.leadOnes, BoolCircuit.CircuitEval.leadOnes_replicate, BoolCircuit.CircuitEval.ofFn_getD_eq_take, BoolCircuit.CircuitEval.verdict, BoolCircuit.DAGCircuit, BoolCircuit.DAGCircuit.IsFaninTwo, BoolCircuit.DAGCircuit.encode, BoolCircuit.DAGCircuit.eval, BoolCircuit.DAGCircuit.values, BoolCircuit.encodeNat, Turing.pairDecode_pairEncode, Turing.pairEncode}
\difficulty{6}
For any Boolean parameter $\mathsf{exact}$ and bit strings $\mathsf{code}$ and $x$, the
verdict on the pairing of $\mathsf{code}$ with $x$ is true if and only if there exist a
natural number $n$ and a directed acyclic Boolean circuit $C$ with $n$ input vertices
such that $C$ is well formed with fan-in at most two, its encoding is $\mathsf{code}$,
$n\le |x|$, and, if $\mathsf{exact}$ is true, then $n=|x|$, while $C$ evaluates to true
on the first $n$ bits of $x$ (taking any out-of-range input lookup to be false).
\end{theorem}
